\begin{proof}
    First it must be shown that for each $f \in P( \Gamma )$, there exists a $g \in P( \Sigma )$ such that $f \approx_E g$.
    The proof follows by structural induction on $P( \Gamma )$.
    \begin{itemize}
    \item \textbf{Identity}.
          Assume that $f = 1_n$.
          Since $P( \Sigma )$ is a subcategory of $P( \Gamma )$, then $f \in P( \Sigma )$.
          Since $( \approx_E )$ is reflexive, then $f \approx_E f$.
    \item \textbf{Symmetries}.
          Assume that $f = \sigma$.
          Since $P( \Sigma )$ is a subcategory of $P( \Gamma )$, then $f \in P( \Sigma )$.
          Since $( \approx_E )$ is reflexive, then $f \approx_E f$.
    \item \textbf{Generators}.
          Assume that $f = x$ for some $x \in \Gamma$.
          Then by assumption, there exists some $g \in P( \Sigma )$ such that $f \approx_E g$.
    \item \textbf{Sequential Composition}.
          Assume that $f = b \circ a$ for $a: n \to s$ and $b: s \to m$.
          Then by the inductive hypothesis, there exists some $g \in P( \Sigma )( n, s )$ and $h \in P( \Sigma )( s, m )$ such that $a \approx_E g$ and $b \approx_E h$.
          Then $f \approx_E h \circ a \approx_E h \circ g$.
          Then $f \approx_E h \circ g$ by the transitivity of $( \approx_E )$.
          Moreover, $h \circ g \in P( \Sigma )( n, m )$.
    \item \textbf{Parallel Composition}.
          Assume that $f = a \boxtimes b$ for $a: n \to m$ and $b: s \to t$.
          Then by the inductive hypothesis, there exists some $g \in P( \Sigma )( n, m )$ and $h \in P( \Sigma )( s, t )$ such that $a \approx_E g$ and $b \approx_E h$.
          Then $f \approx_E g \otimes b \approx_E g \otimes h$.
          Then $f \approx_E g \otimes h$ by the transitivity of $( \approx_E )$.
          Moreover, $g \otimes h \in P( \Sigma )( n + s, m + t )$.
    \end{itemize}
    Then for all $f \in P( \Gamma )$, there exists a $g \in P( \Sigma )$ such that $f \approx_E g$.
    It remains to be shown that $( \Gamma, E )$ is complete with respect to $\sem{-}$.
    Let $f \in P( \Gamma )( n, m )$ and $g \in P( \Gamma )( n, m )$ with $\sem{f} = \sem{g}$.
    By the first result, there exists some $h \in P( \Sigma )( n, m )$ and $k \in P( \Gamma )( n, m )$ such that $f \approx_E h$ and $g \approx_E k$.
    Since $E$ is sound with respect to $E$, then $\sem{f} =\sem{h}$ and $\sem{g} = \sem{k}$.
    Then $\sem{h} = \sem{f} = \sem{g} = \sem{k}$.
    Since $( \Sigma, E )$ is complete with respect to $\sem{-}$, then $h \approx_E k$.
    Since $( \approx_E )$ is transitive, then $f \approx_E g$.
    Since $f$ and $g$ were arbitrary, then $( \Gamma, E )$ is complete with respect to $\sem{-}$.
\end{proof}
